\documentclass{ximera}
%geschreven 19/11/26
\title{Herhaling: vraagstukken over meerdere hoofdstukken}
\author{Daan Lipkens}

\begin{document}
\begin{abstract}
    Vraagstukken die leerstof van de elektrische krachtwerking, het elektrisch veld en de potentiaal combineren.
\end{abstract}
\maketitle

%———————————————————————————————————————
% GEÏNTEGREERDE HERHALINGSOEFENINGEN H1-H3
% Elke oefening combineert minstens twee hoofdstukken (zie commentaar).
% Van makkelijk naar moeilijk.
%———————————————————————————————————————
\section{Vraagstukken over meerdere hoofdstukken}

In elk vraagstuk hieronder heb je leerstof uit meer dan één hoofdstuk nodig. De deelvragen bouwen op elkaar verder: gebruik je resultaat van een vorige deelvraag in de volgende.

%———————————————————————————————————————
% GEÏNTEGREERD 1 (H1 + H2): opladen via contact, coulombkracht, veld van twee gelijke ladingen, geleider
\begin{exercise}
    Twee identieke metalen bolletjes $A$ en $B$ hangen aan een isolerend draadje. Bolletje $A$ heeft een lading van $+8{,}0 \, \mathrm{nC}$, bolletje $B$ een lading van $-2{,}0 \, \mathrm{nC}$. Je laat de bolletjes elkaar even raken en brengt ze daarna op $20 \, \mathrm{cm}$ van elkaar. Beschouw de bolletjes als puntladingen.

    \begin{question}
        Bereken de lading van elk bolletje na het contact.

        \[ Q_A = Q_B = \answer[tolerance=0.08,onlineshowanswerbutton]{3.0} \cdot 10^{-9} \, \mathrm{C} \]

        \begin{hint}
            Hoe groot is de totale lading van de twee bolletjes samen, voor en na het contact?
        \end{hint}

        \begin{solution}\nl
            \underline{Gegevens:}
            \begin{align*}
                & Q_A = +8{,}0 \cdot 10^{-9} \, \mathrm{C} & & Q_B = -2{,}0 \cdot 10^{-9} \, \mathrm{C}
            \end{align*}

            \underline{Gevraagd:} $Q_A'$ en $Q_B'$

            \underline{Oplossing:}\\
            Volgens de wet van behoud van lading blijft de totale lading gelijk:
            \begin{align*}
                Q_{\text{tot}} = Q_A + Q_B = +8{,}0 \cdot 10^{-9} \, \mathrm{C} + (-2{,}0 \cdot 10^{-9} \, \mathrm{C}) = +6{,}0 \cdot 10^{-9} \, \mathrm{C}
            \end{align*}
            De bolletjes zijn identieke geleiders, dus de totale lading verdeelt zich gelijk over beide:
            \begin{align*}
                Q_A' = Q_B' = \frac{Q_{\text{tot}}}{2} = \frac{+6{,}0 \cdot 10^{-9} \, \mathrm{C}}{2} = \boxed{+3{,}0 \cdot 10^{-9} \, \mathrm{C}}
            \end{align*}
        \end{solution}
    \end{question}

    \begin{question}
        Hoeveel elektronen zijn er tijdens het contact overgegaan en in welke zin?

        \[ n = \answer[tolerance=0.08,onlineshowanswerbutton]{3.1} \cdot 10^{10} \]

        De elektronen gaan van \wordChoice{\choice{$A$ naar $B$}\choice[correct]{$B$ naar $A$}}.

        \begin{solution}\nl
            \underline{Gegevens:}
            \begin{align*}
                & Q_B = -2{,}0 \cdot 10^{-9} \, \mathrm{C} & & Q_B' = +3{,}0 \cdot 10^{-9} \, \mathrm{C} & & e = 1{,}60 \cdot 10^{-19} \, \mathrm{C}
            \end{align*}

            \underline{Gevraagd:} $n$

            \underline{Oplossing:}\\
            Enkel elektronen kunnen zich verplaatsen. De lading van $B$ stijgt van $-2{,}0 \, \mathrm{nC}$ naar $+3{,}0 \, \mathrm{nC}$: $B$ heeft dus elektronen afgegeven aan $A$. De verandering van de lading van $B$ is:
            \begin{align*}
                \Delta Q_B = Q_B' - Q_B = +3{,}0 \cdot 10^{-9} \, \mathrm{C} - (-2{,}0 \cdot 10^{-9} \, \mathrm{C}) = +5{,}0 \cdot 10^{-9} \, \mathrm{C}
            \end{align*}
            Elk elektron draagt een lading met grootte $e$:
            \begin{align*}
                n = \frac{|\Delta Q_B|}{e} = \frac{5{,}0 \cdot 10^{-9} \, \mathrm{C}}{1{,}60 \cdot 10^{-19} \, \mathrm{C}} = \boxed{3{,}1 \cdot 10^{10}}
            \end{align*}
        \end{solution}
    \end{question}

    \begin{question}
        Bereken de grootte van de elektrische kracht tussen de bolletjes op $20 \, \mathrm{cm}$ van elkaar.

        \[ F_{AB} = \answer[tolerance=0.05,onlineshowanswerbutton]{2.0} \cdot 10^{-6} \, \mathrm{N} \]

        \begin{solution}\nl
            \underline{Gegevens:}
            \begin{align*}
                & Q_A' = Q_B' = +3{,}0 \cdot 10^{-9} \, \mathrm{C} & & r = 20 \, \mathrm{cm} = 0{,}20 \, \mathrm{m}
            \end{align*}

            \underline{Gevraagd:} $F_{AB}$

            \underline{Oplossing:}
            \begin{align*}
                F_{AB} &= k \cdot \frac{|Q_A'| \cdot |Q_B'|}{r^2} \\
                &= 8{,}99 \cdot 10^{9} \, \mathrm{\frac{N \cdot m^2}{C^2}} \cdot \frac{3{,}0 \cdot 10^{-9} \, \mathrm{C} \cdot 3{,}0 \cdot 10^{-9} \, \mathrm{C}}{(0{,}20 \, \mathrm{m})^2} = \boxed{2{,}0 \cdot 10^{-6} \, \mathrm{N}}
            \end{align*}
            Beide bolletjes zijn nu positief, dus ze stoten elkaar af.
        \end{solution}
    \end{question}

    \begin{question}
        Hoe groot is het totale elektrisch veld in het midden tussen de twee bolletjes?

        \begin{multipleChoice}
            \choice[correct]{Nul.}
            \choice{Twee keer het veld van één bolletje.}
            \choice{Even groot als het veld van één bolletje.}
        \end{multipleChoice}

        \begin{solution}\nl
            In het midden liggen beide bolletjes op dezelfde afstand en hebben ze dezelfde lading. Hun velden zijn dus even groot. Beide ladingen zijn positief, dus elk veld wijst weg van zijn eigen bolletje: de velden zijn tegengesteld en heffen elkaar op.
        \end{solution}
    \end{question}

    \begin{question}
        Hoe groot is het elektrisch veld binnenin het metaal van bolletje $A$?

        \begin{multipleChoice}
            \choice[correct]{Nul, want $A$ is een geleider in elektrostatisch evenwicht.}
            \choice{Het grootst in het middelpunt van $A$.}
            \choice{Even groot als net buiten het oppervlak van $A$.}
        \end{multipleChoice}

        \begin{solution}\nl
            In een geleider in elektrostatisch evenwicht is het veld binnenin steeds nul. Was dat niet zo, dan zouden de vrije elektronen een kracht ondervinden en zich verder verplaatsen. De lading van het bolletje zit aan het oppervlak.
        \end{solution}
    \end{question}
\end{exercise}
%———————————————————————————————————————

%———————————————————————————————————————
% GEÏNTEGREERD 2 (H1 + H2): coulombkracht, veld als kracht per lading, nulpunt van het veld
\begin{exercise}
    Twee puntladingen $Q_1 = +6{,}0 \, \mathrm{nC}$ en $Q_2 = -2{,}0 \, \mathrm{nC}$ liggen op $12 \, \mathrm{cm}$ van elkaar op een rechte lijn.

    \begin{question}
        Bereken de grootte van de kracht die $Q_1$ op $Q_2$ uitoefent.

        \[ F_{12} = \answer[tolerance=0.2,onlineshowanswerbutton]{7.5} \cdot 10^{-6} \, \mathrm{N} \]

        \begin{solution}\nl
            \underline{Gegevens:}
            \begin{align*}
                & Q_1 = +6{,}0 \cdot 10^{-9} \, \mathrm{C} & & Q_2 = -2{,}0 \cdot 10^{-9} \, \mathrm{C} & & r_{12} = 12 \, \mathrm{cm} = 0{,}12 \, \mathrm{m}
            \end{align*}

            \underline{Gevraagd:} $F_{12}$

            \underline{Oplossing:}
            \begin{align*}
                F_{12} &= k \cdot \frac{|Q_1| \cdot |Q_2|}{r_{12}^2} \\
                &= 8{,}99 \cdot 10^{9} \, \mathrm{\frac{N \cdot m^2}{C^2}} \cdot \frac{6{,}0 \cdot 10^{-9} \, \mathrm{C} \cdot 2{,}0 \cdot 10^{-9} \, \mathrm{C}}{(0{,}12 \, \mathrm{m})^2} = \boxed{7{,}5 \cdot 10^{-6} \, \mathrm{N}}
            \end{align*}
        \end{solution}
    \end{question}

    \begin{question}
        Bereken de grootte van het elektrisch veld van $Q_1$ op de plaats van $Q_2$.

        \[ E_1 = \answer[tolerance=0.09,onlineshowanswerbutton]{3.7} \cdot 10^{3} \, \mathrm{N/C} \]

        \begin{hint}
            Er zijn twee manieren: met de kracht uit de vorige deelvraag of met de formule voor het veld van een puntlading. Welke lading is hier de bronlading en welke de testlading?
        \end{hint}

        \begin{solution}\nl
            \underline{Gegevens:}
            \begin{align*}
                & F_{12} = 7{,}49 \cdot 10^{-6} \, \mathrm{N} & & Q_t = Q_2 = -2{,}0 \cdot 10^{-9} \, \mathrm{C}
            \end{align*}

            \underline{Gevraagd:} $E_1$ op de plaats van $Q_2$

            \underline{Oplossing:}\\
            $Q_1$ is de bronlading en $Q_2$ de testlading. Het veld is de kracht per eenheid van lading:
            \begin{align*}
                E_1 = \frac{F_{12}}{|Q_2|} = \frac{7{,}49 \cdot 10^{-6} \, \mathrm{N}}{2{,}0 \cdot 10^{-9} \, \mathrm{C}} = \boxed{3{,}7 \cdot 10^{3} \, \mathrm{N/C}}
            \end{align*}
            Controle met het veld van een puntlading:
            \begin{align*}
                E_1 = k \frac{|Q_1|}{r^2} = \frac{8{,}99 \cdot 10^{9} \, \mathrm{\frac{N \cdot m^2}{C^2}} \cdot 6{,}0 \cdot 10^{-9} \, \mathrm{C}}{(0{,}12 \, \mathrm{m})^2} = 3{,}7 \cdot 10^{3} \, \mathrm{N/C}
            \end{align*}
        \end{solution}
    \end{question}

    \begin{question}
        Het veld van $Q_1$ op de plaats van $Q_2$ en de kracht op $Q_2$ hebben
        \wordChoice{\choice{dezelfde zin}\choice[correct]{een tegengestelde zin}}.

        \begin{solution}\nl
            $Q_2$ is negatief. Op een negatieve lading werkt een kracht tegengesteld aan het veld. Het veld van $Q_1$ wijst weg van $Q_1$, de kracht op $Q_2$ wijst naar $Q_1$ toe: de ladingen trekken elkaar aan.
        \end{solution}
    \end{question}

    \begin{question}
        Op de lijn door de twee ladingen ligt één punt waar het totale elektrisch veld nul is. Waar ligt dat punt?

        \begin{multipleChoice}
            \choice{Tussen de twee ladingen.}
            \choice{Aan de buitenkant, aan de kant van $Q_1$.}
            \choice[correct]{Aan de buitenkant, aan de kant van $Q_2$.}
        \end{multipleChoice}

        \begin{solution}\nl
            Tussen twee ladingen met een tegengesteld teken wijzen beide velden in dezelfde zin: ze kunnen elkaar daar niet opheffen. Buiten de ladingen zijn de velden wel tegengesteld. Om even groot te zijn, moet het punt dichter bij de kleinste lading liggen, dus aan de kant van $Q_2$.
        \end{solution}
    \end{question}

    \begin{question}
        Bereken de afstand $x$ van dat punt tot $Q_2$.

        \[ x = \answer[tolerance=0.004,onlineshowanswerbutton]{0.16} \, \mathrm{m} \]

        \begin{hint}
            Als het punt op een afstand $x$ van $Q_2$ ligt, hoe ver ligt het dan van $Q_1$?
        \end{hint}

        \begin{solution}\nl
            \underline{Gegevens:}
            \begin{align*}
                & Q_1 = +6{,}0 \cdot 10^{-9} \, \mathrm{C} & & Q_2 = -2{,}0 \cdot 10^{-9} \, \mathrm{C} & & d = 0{,}12 \, \mathrm{m}
            \end{align*}

            \underline{Gevraagd:} $x$, de afstand tot $Q_2$ waar $E_{\text{tot}} = 0$

            \underline{Oplossing:}\\
            Het punt ligt op een afstand $x$ van $Q_2$ en op een afstand $d + x$ van $Q_1$. Daar zijn de twee velden even groot:
            \begin{align*}
                k \frac{|Q_1|}{(d+x)^2} &= k \frac{|Q_2|}{x^2} \\
                \left(\frac{d+x}{x}\right)^2 &= \frac{|Q_1|}{|Q_2|} = \frac{6{,}0 \cdot 10^{-9} \, \mathrm{C}}{2{,}0 \cdot 10^{-9} \, \mathrm{C}} = 3{,}0 \\
                \frac{d+x}{x} &= \sqrt{3{,}0} = 1{,}73
            \end{align*}
            Hieruit volgt $d + x = 1{,}73 \, x$, dus $d = 0{,}73 \, x$ en
            \begin{align*}
                x = \frac{d}{\sqrt{3{,}0} - 1} = \frac{0{,}12 \, \mathrm{m}}{0{,}73} = \boxed{1{,}6 \cdot 10^{-1} \, \mathrm{m}}
            \end{align*}
        \end{solution}
    \end{question}
\end{exercise}
%———————————————————————————————————————

%———————————————————————————————————————
% GEÏNTEGREERD 3 (H2 + H3): homogeen veld met figuur, kracht, potentiaal, spanning en arbeid
\begin{exercise}
    Twee verticale evenwijdige platen staan op $10 \, \mathrm{cm}$ van elkaar. De linkerplaat is positief, de rechterplaat negatief. Tussen de platen staat een spanning van $50 \, \mathrm{V}$. De punten $A$, $B$ en $C$ liggen tussen de platen, zoals op de figuur hieronder. Het referentiepunt ligt op de negatieve plaat. In $A$ bevindt zich een lading $Q = -4{,}0 \, \mathrm{nC}$.

    \begin{center}
            \begin{tikzpicture}
                \draw[thick] (0,-1.6) -- (0,1.6);
                \draw[thick] (5,-1.6) -- (5,1.6);
                \foreach \y in {-1.2,-0.6,0,0.6,1.2} {
                    \node[left] at (-0.1,\y) {$+$};
                    \node[right] at (5.1,\y) {$-$};
                    \draw[Eveldlijn] (0,\y) -- (5,\y);
                }
                \fill (1,0.9) circle (1.5pt) node[right=2pt] {$A$};
                \fill (3.5,0.9) circle (1.5pt) node[right=2pt] {$B$};
                \fill (3.5,-0.9) circle (1.5pt) node[right=2pt] {$C$};
                \draw[<->] (0,-2.0) -- (1,-2.0) node[midway,below] {$2{,}0\,\mathrm{cm}$};
                \draw[<->] (0,-2.7) -- (3.5,-2.7) node[midway,below] {$7{,}0\,\mathrm{cm}$};
                \draw[<->] (0,-3.4) -- (5,-3.4) node[midway,below] {$10\,\mathrm{cm}$};
                \draw[dashed, gray] (1,0.9) -- (1,-2.0);
                \draw[dashed, gray] (3.5,-0.9) -- (3.5,-2.7);
            \end{tikzpicture}
    \end{center}
    De afstanden op de figuur zijn gemeten vanaf de positieve plaat.

    \begin{question}
        Bereken de grootte van het elektrisch veld tussen de platen.

        \[ E = \answer[tolerance=10,onlineshowanswerbutton]{500} \, \mathrm{N/C} \]

        \begin{solution}\nl
            \underline{Gegevens:}
            \begin{align*}
                & U = 50 \, \mathrm{V} & & \ell = 10 \, \mathrm{cm} = 0{,}10 \, \mathrm{m}
            \end{align*}

            \underline{Gevraagd:} $E$

            \underline{Oplossing:}
            \begin{align*}
                E = \frac{U}{\ell} = \frac{50 \, \mathrm{V}}{0{,}10 \, \mathrm{m}} = \boxed{5{,}0 \cdot 10^{2} \, \mathrm{N/C}}
            \end{align*}
        \end{solution}
    \end{question}

    \begin{question}
        Bereken de grootte van de elektrische kracht op de lading in $A$.

        \[ F_e = \answer[tolerance=0.05,onlineshowanswerbutton]{2.0} \cdot 10^{-6} \, \mathrm{N} \]

        De kracht wijst naar de \wordChoice{\choice[correct]{positieve plaat}\choice{negatieve plaat}}.

        \begin{solution}\nl
            \underline{Gegevens:}
            \begin{align*}
                & Q = -4{,}0 \cdot 10^{-9} \, \mathrm{C} & & E = 5{,}0 \cdot 10^{2} \, \mathrm{N/C}
            \end{align*}

            \underline{Gevraagd:} $F_e$

            \underline{Oplossing:}
            \begin{align*}
                F_e = |Q| \cdot E = 4{,}0 \cdot 10^{-9} \, \mathrm{C} \cdot 5{,}0 \cdot 10^{2} \, \frac{\mathrm{N}}{\mathrm{C}} = \boxed{2{,}0 \cdot 10^{-6} \, \mathrm{N}}
            \end{align*}
            De lading is negatief, dus de kracht is tegengesteld aan het veld: naar links, naar de positieve plaat.
        \end{solution}
    \end{question}

    \begin{question}
        Bereken de potentiaal in $A$ en in $B$.

        \[ V_A = \answer[tolerance=1,onlineshowanswerbutton]{40} \, \mathrm{J/C} \qquad V_B = \answer[tolerance=0.4,onlineshowanswerbutton]{15} \, \mathrm{J/C} \]

        \begin{hint}
            De afstanden op de figuur zijn gemeten vanaf de positieve plaat. Ten opzichte van welke plaat heb je de afstand nodig?
        \end{hint}

        \begin{solution}\nl
            \underline{Gegevens:}
            \begin{align*}
                & E = 5{,}0 \cdot 10^{2} \, \mathrm{N/C} \\
                & d_A = 10 \, \mathrm{cm} - 2{,}0 \, \mathrm{cm} = 0{,}080 \, \mathrm{m} & & d_B = 10 \, \mathrm{cm} - 7{,}0 \, \mathrm{cm} = 0{,}030 \, \mathrm{m}
            \end{align*}

            \underline{Gevraagd:} $V_A$ en $V_B$

            \underline{Oplossing:}\\
            De potentiaal hangt af van de afstand tot de negatieve plaat (het referentiepunt):
            \begin{align*}
                V_A &= E \cdot d_A = 5{,}0 \cdot 10^{2} \, \frac{\mathrm{N}}{\mathrm{C}} \cdot 0{,}080 \, \mathrm{m} = \boxed{4{,}0 \cdot 10^{1} \, \mathrm{J/C}} \\
                V_B &= E \cdot d_B = 5{,}0 \cdot 10^{2} \, \frac{\mathrm{N}}{\mathrm{C}} \cdot 0{,}030 \, \mathrm{m} = \boxed{1{,}5 \cdot 10^{1} \, \mathrm{J/C}}
            \end{align*}
        \end{solution}
    \end{question}

    \begin{question}
        De lading beweegt van $A$ naar $B$. Bereken de arbeid die het elektrisch veld daarbij verricht.

        \[ W_{A \to B} = \answer[tolerance=0.03,onlineshowanswerbutton]{-1.0} \cdot 10^{-7} \, \mathrm{J} \]

        \begin{solution}\nl
            \underline{Gegevens:}
            \begin{align*}
                & Q = -4{,}0 \cdot 10^{-9} \, \mathrm{C} & & V_A = 40 \, \mathrm{J/C} & & V_B = 15 \, \mathrm{J/C}
            \end{align*}

            \underline{Gevraagd:} $W_{A \to B}$

            \underline{Oplossing:}\\
            Eerst de spanning tussen $A$ en $B$:
            \begin{align*}
                U_{AB} = V_A - V_B = 40 \, \mathrm{J/C} - 15 \, \mathrm{J/C} = 25 \, \mathrm{V}
            \end{align*}
            De arbeid van het veld:
            \begin{align*}
                W_{A \to B} = Q \cdot U_{AB} = -4{,}0 \cdot 10^{-9} \, \mathrm{C} \cdot 25 \, \mathrm{V} = \boxed{-1{,}0 \cdot 10^{-7} \, \mathrm{J}}
            \end{align*}
            De arbeid is negatief: de lading beweegt van $A$ naar $B$ tegen de elektrische kracht in (die wijst naar de positieve plaat). Er is dus een externe kracht nodig. Dat klopt met de vorige deelvraag over de zin van de kracht.
        \end{solution}
    \end{question}

    \begin{question}
        Daarna beweegt de lading van $B$ naar $C$. De arbeid die het veld daarbij verricht, is
        \wordChoice{\choice{positief}\choice{negatief}\choice[correct]{nul}}.

        \begin{solution}\nl
            $B$ en $C$ liggen op dezelfde afstand van de negatieve plaat, dus op hetzelfde equipotentiaaloppervlak. De spanning $U_{BC}$ is nul en dus ook de arbeid $W = Q \cdot U_{BC}$. De verplaatsing staat loodrecht op de elektrische kracht.
        \end{solution}
    \end{question}
\end{exercise}
%———————————————————————————————————————

%———————————————————————————————————————
% GEÏNTEGREERD 4 (H1 + H2 + H3): zwevend stofdeeltje, aantal elektronen gegeven, spanning gevraagd
\begin{exercise}
    Een klein stofdeeltje met massa $4{,}9 \cdot 10^{-15} \, \mathrm{kg}$ heeft $6$ elektronen te veel. Het bevindt zich tussen twee horizontale evenwijdige platen op $8{,}0 \, \mathrm{mm}$ van elkaar. Bij een bepaalde spanning tussen de platen blijft het deeltje zweven. Gebruik $g = 9{,}81 \, \mathrm{N/kg}$.

    \begin{question}
        Bereken de lading van het stofdeeltje.

        \[ Q = \answer[tolerance=0.2,onlineshowanswerbutton]{-9.6} \cdot 10^{-19} \, \mathrm{C} \]

        \begin{solution}\nl
            \underline{Gegevens:}
            \begin{align*}
                & n = 6 & & e = 1{,}60 \cdot 10^{-19} \, \mathrm{C}
            \end{align*}

            \underline{Gevraagd:} $Q$

            \underline{Oplossing:}
            \begin{align*}
                Q = -n \cdot e = -6 \cdot 1{,}60 \cdot 10^{-19} \, \mathrm{C} = \boxed{-9{,}6 \cdot 10^{-19} \, \mathrm{C}}
            \end{align*}
        \end{solution}
    \end{question}

    \begin{question}
        Welke plaat moet positief zijn opdat het deeltje kan zweven?

        \wordChoice{\choice[correct]{de bovenste plaat}\choice{de onderste plaat}}

        \begin{hint}
            In welke zin moet de elektrische kracht werken? Welke zin heeft het veld dan, voor een negatieve lading?
        \end{hint}

        \begin{solution}\nl
            De zwaartekracht wijst naar beneden, dus de elektrische kracht moet naar boven wijzen. Het deeltje is negatief: de kracht is tegengesteld aan het veld. Het veld wijst dus naar beneden, van de positieve naar de negatieve plaat. De bovenste plaat is positief.
        \end{solution}
    \end{question}

    \begin{question}
        Bereken de grootte van het elektrisch veld tussen de platen.

        \[ E = \answer[tolerance=0.1,onlineshowanswerbutton]{5.0} \cdot 10^{4} \, \mathrm{N/C} \]

        \begin{hint}
            Het deeltje zweeft. Wat weet je dan over de krachten die op het deeltje werken?
        \end{hint}

        \begin{solution}\nl
            \underline{Gegevens:}
            \begin{align*}
                & m = 4{,}9 \cdot 10^{-15} \, \mathrm{kg} & & g = 9{,}81 \, \mathrm{N/kg} & & |Q| = 9{,}6 \cdot 10^{-19} \, \mathrm{C}
            \end{align*}

            \underline{Gevraagd:} $E$

            \underline{Oplossing:}\\
            Het deeltje zweeft, dus de resulterende kracht is nul. De elektrische kracht en de zwaartekracht zijn even groot:
            \begin{align*}
                F_e &= F_z \\
                |Q| \cdot E &= m \cdot g \\
                E &= \frac{m \cdot g}{|Q|} = \frac{4{,}9 \cdot 10^{-15} \, \mathrm{kg} \cdot 9{,}81 \, \mathrm{N/kg}}{9{,}6 \cdot 10^{-19} \, \mathrm{C}} = \boxed{5{,}0 \cdot 10^{4} \, \mathrm{N/C}}
            \end{align*}
        \end{solution}
    \end{question}

    \begin{question}
        Bereken de spanning tussen de platen.

        \[ U = \answer[tolerance=10,onlineshowanswerbutton]{400} \, \mathrm{V} \]

        \begin{solution}\nl
            \underline{Gegevens:}
            \begin{align*}
                & E = 5{,}01 \cdot 10^{4} \, \mathrm{N/C} & & \ell = 8{,}0 \, \mathrm{mm} = 8{,}0 \cdot 10^{-3} \, \mathrm{m}
            \end{align*}

            \underline{Gevraagd:} $U$

            \underline{Oplossing:}
            \begin{align*}
                U = E \cdot \ell = 5{,}01 \cdot 10^{4} \, \frac{\mathrm{N}}{\mathrm{C}} \cdot 8{,}0 \cdot 10^{-3} \, \mathrm{m} = \boxed{4{,}0 \cdot 10^{2} \, \mathrm{V}}
            \end{align*}
        \end{solution}
    \end{question}
\end{exercise}
%———————————————————————————————————————

%———————————————————————————————————————
% GEÏNTEGREERD 5 (H2 + H3, 2 uur): proton tussen platen, energieomzetting, vergelijking met alfadeeltje
\begin{exercise}
    Tussen twee evenwijdige platen op $5{,}0 \, \mathrm{cm}$ van elkaar heerst een homogeen elektrisch veld met grootte $4{,}0 \, \mathrm{kN/C}$. Een proton vertrekt vanuit rust aan de positieve plaat. De massa van een proton is $1{,}67 \cdot 10^{-27} \, \mathrm{kg}$. Enkel de elektrische kracht werkt op het proton.

    \begin{question}
        Bereken de spanning tussen de platen.

        \[ U = \answer[tolerance=5,onlineshowanswerbutton]{200} \, \mathrm{V} \]

        \begin{solution}\nl
            \underline{Gegevens:}
            \begin{align*}
                & E = 4{,}0 \, \mathrm{kN/C} = 4{,}0 \cdot 10^{3} \, \mathrm{N/C} & & \ell = 5{,}0 \, \mathrm{cm} = 5{,}0 \cdot 10^{-2} \, \mathrm{m}
            \end{align*}

            \underline{Gevraagd:} $U$

            \underline{Oplossing:}
            \begin{align*}
                U = E \cdot \ell = 4{,}0 \cdot 10^{3} \, \frac{\mathrm{N}}{\mathrm{C}} \cdot 5{,}0 \cdot 10^{-2} \, \mathrm{m} = \boxed{2{,}0 \cdot 10^{2} \, \mathrm{V}}
            \end{align*}
        \end{solution}
    \end{question}

    \begin{question}
        Bereken de potentiële elektrische energie van het proton bij de start. Het referentiepunt ligt op de negatieve plaat.

        \[ E_{\text{pot}} = \answer[tolerance=0.08,onlineshowanswerbutton]{3.2} \cdot 10^{-17} \, \mathrm{J} \]

        \begin{solution}\nl
            \underline{Gegevens:}
            \begin{align*}
                & Q = +e = +1{,}60 \cdot 10^{-19} \, \mathrm{C} & & E = 4{,}0 \cdot 10^{3} \, \mathrm{N/C} & & d = 5{,}0 \cdot 10^{-2} \, \mathrm{m}
            \end{align*}

            \underline{Gevraagd:} $E_{\text{pot}}$

            \underline{Oplossing:}\\
            Het proton start aan de positieve plaat, dus op de volledige plaatafstand van de negatieve plaat:
            \begin{align*}
                E_{\text{pot}} &= Q \cdot E \cdot d \\
                &= 1{,}60 \cdot 10^{-19} \, \mathrm{C} \cdot 4{,}0 \cdot 10^{3} \, \frac{\mathrm{N}}{\mathrm{C}} \cdot 5{,}0 \cdot 10^{-2} \, \mathrm{m} = \boxed{3{,}2 \cdot 10^{-17} \, \mathrm{J}}
            \end{align*}
        \end{solution}
    \end{question}

    \begin{question}
        Bereken de snelheid waarmee het proton de negatieve plaat bereikt.

        \[ v = \answer[tolerance=0.05,onlineshowanswerbutton]{2.0} \cdot 10^{5} \, \mathrm{m/s} \]

        \begin{hint}
            Hoe groot is de potentiële energie van het proton als het aan de negatieve plaat aankomt? Waar is die energie naartoe?
        \end{hint}

        \begin{solution}\nl
            \underline{Gegevens:}
            \begin{align*}
                & E_{\text{pot,begin}} = 3{,}2 \cdot 10^{-17} \, \mathrm{J} & & m_p = 1{,}67 \cdot 10^{-27} \, \mathrm{kg} & & v_{\text{begin}} = 0 \, \mathrm{m/s}
            \end{align*}

            \underline{Gevraagd:} $v$ aan de negatieve plaat

            \underline{Oplossing:}\\
            Aan de negatieve plaat is de potentiële energie nul. Enkel de elektrische kracht werkt, dus de volledige potentiële energie is omgezet in kinetische energie: $E_{\text{kin}} = 3{,}2 \cdot 10^{-17} \, \mathrm{J}$. Met $E_{\text{kin}} = \frac{m \cdot v^2}{2}$:
            \begin{align*}
                v = \sqrt{\frac{2 \cdot E_{\text{kin}}}{m_p}} = \sqrt{\frac{2 \cdot 3{,}2 \cdot 10^{-17} \, \mathrm{J}}{1{,}67 \cdot 10^{-27} \, \mathrm{kg}}} = \boxed{2{,}0 \cdot 10^{5} \, \mathrm{m/s}}
            \end{align*}
        \end{solution}
    \end{question}

    \begin{question}
        Je vervangt het proton door een alfadeeltje. Dat heeft een twee keer zo grote lading en een ongeveer vier keer zo grote massa. Het vertrekt ook vanuit rust aan de positieve plaat. Wat geldt voor het alfadeeltje aan de negatieve plaat, vergeleken met het proton?

        \begin{multipleChoice}
            \choice{Dezelfde kinetische energie en dezelfde snelheid.}
            \choice{Twee keer zoveel kinetische energie en een grotere snelheid.}
            \choice[correct]{Twee keer zoveel kinetische energie en een kleinere snelheid.}
            \choice{De helft van de kinetische energie en een kleinere snelheid.}
        \end{multipleChoice}

        \begin{hint}
            Vergelijk eerst de arbeid van het veld, die hangt enkel af van de lading en de spanning. Gebruik daarna de massa.
        \end{hint}

        \begin{solution}\nl
            De arbeid van het veld is $W = Q \cdot U$. De spanning is dezelfde en de lading is twee keer zo groot: het alfadeeltje krijgt twee keer zoveel kinetische energie. Uit $E_{\text{kin}} = \frac{m \cdot v^2}{2}$ volgt $v^2 = \frac{2 \cdot E_{\text{kin}}}{m}$. De teller wordt $2$ keer zo groot, de noemer $4$ keer: $v^2$ wordt $2$ keer zo klein en $v$ dus $\sqrt{2}$ keer zo klein. Het alfadeeltje is trager dan het proton.
        \end{solution}
    \end{question}
\end{exercise}
%———————————————————————————————————————

%———————————————————————————————————————
% GEÏNTEGREERD 6 (H1 + H2 + H3, 2 uur): elektron tussen platen, zwaartekracht verwaarlozen, snelheid, invloed plaatafstand
\begin{exercise}
    Tussen twee evenwijdige platen op $2{,}0 \, \mathrm{cm}$ van elkaar staat een spanning van $300 \, \mathrm{V}$. Een elektron vertrekt vanuit rust aan de negatieve plaat. De massa van een elektron is $9{,}11 \cdot 10^{-31} \, \mathrm{kg}$. Gebruik $g = 9{,}81 \, \mathrm{N/kg}$.

    \begin{question}
        Bereken de grootte van het elektrisch veld tussen de platen.

        \[ E = \answer[tolerance=0.04,onlineshowanswerbutton]{1.5} \cdot 10^{4} \, \mathrm{N/C} \]

        \begin{solution}\nl
            \underline{Gegevens:}
            \begin{align*}
                & U = 300 \, \mathrm{V} & & \ell = 2{,}0 \, \mathrm{cm} = 2{,}0 \cdot 10^{-2} \, \mathrm{m}
            \end{align*}

            \underline{Gevraagd:} $E$

            \underline{Oplossing:}
            \begin{align*}
                E = \frac{U}{\ell} = \frac{300 \, \mathrm{V}}{2{,}0 \cdot 10^{-2} \, \mathrm{m}} = \boxed{1{,}5 \cdot 10^{4} \, \mathrm{N/C}}
            \end{align*}
        \end{solution}
    \end{question}

    \begin{question}
        Bereken de grootte van de elektrische kracht op het elektron.

        \[ F_e = \answer[tolerance=0.06,onlineshowanswerbutton]{2.4} \cdot 10^{-15} \, \mathrm{N} \]

        \begin{solution}\nl
            \underline{Gegevens:}
            \begin{align*}
                & |Q| = e = 1{,}60 \cdot 10^{-19} \, \mathrm{C} & & E = 1{,}5 \cdot 10^{4} \, \mathrm{N/C}
            \end{align*}

            \underline{Gevraagd:} $F_e$

            \underline{Oplossing:}
            \begin{align*}
                F_e = |Q| \cdot E = 1{,}60 \cdot 10^{-19} \, \mathrm{C} \cdot 1{,}5 \cdot 10^{4} \, \frac{\mathrm{N}}{\mathrm{C}} = \boxed{2{,}4 \cdot 10^{-15} \, \mathrm{N}}
            \end{align*}
        \end{solution}
    \end{question}

    \begin{question}
        Hoeveel keer is de elektrische kracht groter dan de zwaartekracht op het elektron?

        \[ \frac{F_e}{F_z} = \answer[tolerance=0.07,onlineshowanswerbutton]{2.7} \cdot 10^{14} \]

        \begin{solution}\nl
            \underline{Gegevens:}
            \begin{align*}
                & F_e = 2{,}4 \cdot 10^{-15} \, \mathrm{N} & & m_e = 9{,}11 \cdot 10^{-31} \, \mathrm{kg} & & g = 9{,}81 \, \mathrm{N/kg}
            \end{align*}

            \underline{Gevraagd:} $\frac{F_e}{F_z}$

            \underline{Oplossing:}
            \begin{align*}
                F_z &= m_e \cdot g = 9{,}11 \cdot 10^{-31} \, \mathrm{kg} \cdot 9{,}81 \, \mathrm{N/kg} = 8{,}94 \cdot 10^{-30} \, \mathrm{N} \\
                \frac{F_e}{F_z} &= \frac{2{,}4 \cdot 10^{-15} \, \mathrm{N}}{8{,}94 \cdot 10^{-30} \, \mathrm{N}} = \boxed{2{,}7 \cdot 10^{14}}
            \end{align*}
            De zwaartekracht op het elektron is verwaarloosbaar klein ten opzichte van de elektrische kracht. Daarom mogen we bij geladen deeltjes in een elektrisch veld de zwaartekracht meestal weglaten.
        \end{solution}
    \end{question}

    \begin{question}
        Bereken de snelheid waarmee het elektron de positieve plaat bereikt.

        \[ v = \answer[tolerance=0.03,onlineshowanswerbutton]{1.0} \cdot 10^{7} \, \mathrm{m/s} \]

        \begin{hint}
            Je kunt de arbeid van het veld berekenen zonder de kracht of de afstand te gebruiken.
        \end{hint}

        \begin{solution}\nl
            \underline{Gegevens:}
            \begin{align*}
                & Q = -e = -1{,}60 \cdot 10^{-19} \, \mathrm{C} & & U_{\text{neg} \to \text{pos}} = -300 \, \mathrm{V} & & m_e = 9{,}11 \cdot 10^{-31} \, \mathrm{kg}
            \end{align*}

            \underline{Gevraagd:} $v$ aan de positieve plaat

            \underline{Oplossing:}\\
            Het elektron beweegt van de negatieve plaat (lage potentiaal) naar de positieve plaat (hoge potentiaal). De spanning van begin- naar eindpunt is dus $-300 \, \mathrm{V}$. De arbeid van het veld:
            \begin{align*}
                W = Q \cdot U = (-1{,}60 \cdot 10^{-19} \, \mathrm{C}) \cdot (-300 \, \mathrm{V}) = 4{,}80 \cdot 10^{-17} \, \mathrm{J}
            \end{align*}
            De zwaartekracht verwaarlozen we (vorige deelvraag). De arbeid wordt volledig omgezet in kinetische energie:
            \begin{align*}
                v = \sqrt{\frac{2 \cdot E_{\text{kin}}}{m_e}} = \sqrt{\frac{2 \cdot 4{,}80 \cdot 10^{-17} \, \mathrm{J}}{9{,}11 \cdot 10^{-31} \, \mathrm{kg}}} = \boxed{1{,}0 \cdot 10^{7} \, \mathrm{m/s}}
            \end{align*}
        \end{solution}
    \end{question}

    \begin{question}
        Je zet de platen twee keer zo ver uit elkaar, bij dezelfde spanning van $300 \, \mathrm{V}$. Wat gebeurt er met de snelheid waarmee het elektron de positieve plaat bereikt?

        \begin{multipleChoice}
            \choice{Ze wordt twee keer zo groot, want het elektron wordt over een langere afstand versneld.}
            \choice{Ze wordt kleiner, want het veld wordt twee keer zo klein.}
            \choice[correct]{Ze blijft gelijk.}
        \end{multipleChoice}

        \begin{solution}\nl
            Het veld wordt wel twee keer zo klein ($E = U/\ell$), dus ook de kracht. Maar die kracht werkt over een twee keer zo grote afstand. De arbeid $W = Q \cdot U$ hangt enkel af van de lading en de spanning. Die zijn niet veranderd. Het elektron krijgt dus dezelfde kinetische energie en bereikt de positieve plaat met dezelfde snelheid.
        \end{solution}
    \end{question}
\end{exercise}
%———————————————————————————————————————

\end{document}
